In the canonical fixed replay, ParkSong-Composite used train-split median
resource/activity processing-time estimates, \texttt{cost\_time\_scale=3600},
\texttt{no\_show\_penalty\_weight=1.0}, \texttt{future\_delay\_weight=0},
\texttt{reservation\_margin=0} and prediction-aware strategic idling. It
created $1{,}580$ reservations, used $474$ of them and ended with no unresolved
reservations, giving a reservation-utilization rate of $30.00\%$. The
prediction-execution audit recorded $2{,}832$ matches and $2{,}213$ mismatches.
The implementation contains a one-step temporal lookahead term, but the
canonical parameterization sets \texttt{future\_delay\_weight=0}; the reported
run should therefore be interpreted as a reservation-based rolling-epoch
approximation, not as a full temporal scheduling or min-cost-flow realization of
Park and Song.

These counters show that the predictive reservation mechanism was exercised, but
they also explain why it did not dominate the reactive methods. With the
RF-optimized leakage-free branching artifact, ParkSong completed all $380/380$
fixed-route instances across the five seeds, with mean cycle time of $32.2304$ h
and mean waiting time of $73.3656$ min. This improves completion relative to the
previous corrected run, but the higher timing metrics show that anticipatory
reservations were costly on this workload. A likely contributor is low realized
reservation usefulness: many candidate reservations were filtered by the
current-task comparison and only about one third of created reservations were
eventually consumed. The processing-time sampler remains a limitation because a
large share of simulator-side processing-time samples still uses fallback
estimates rather than exact resource/activity models.
